% TOP_PROBLEM: {"id":30007575,"problem_number":"LOCAL-30007575","title":"Long-term contracts versus spot markets","category_id":21,"category":{"id":21,"name":"economics","display_name":"Economics","description":"Open economic theory statements, published research agendas, and proposed scoped specifications; question provenance and historical priority are qualified per record.","slug":"economics","order_index":21,"created_at":"2026-10-08T00:00:00Z"},"status":"research_question","external_url":null,"legacy_ids":[],"published":false,"statement":"Identify when offered long-term contracts improve net relationship surplus relative to spot trade and whether measured belief disagreement helps explain limited adoption.","economics":{"original_id":64,"field":"Mechanism design and market design","kind":"empirical","formalization_basis":"adapted_research_question","source_status":"research_agenda","priority_status":"earliest_found","priority_note":"This is the earliest source in which the relevant agenda was verified during this audit; absolute priority has not been established.","earliest_verified_reference":{"authors":["Dirk Bergemann","Juuso Välimäki"],"title":"Dynamic Mechanism Design: An Introduction","year":2018,"url":"https://cowles.yale.edu/sites/default/files/2022-09/d2102-r.pdf","doi":"10.1257/jel.20180892","locator":"§8 Concluding Remarks, pp.49–51 of June19,2018 revision","evidence_summary":"The conclusion explicitly identifies weaker common-prior assumptions, sequential participation, outside resale, and limited adoption of long-term contracts as research questions."},"current_reference":{"authors":["Dirk Bergemann","Juuso Välimäki"],"title":"Dynamic Mechanism Design: An Introduction","year":2018,"url":"https://cowles.yale.edu/sites/default/files/2022-09/d2102-r.pdf","doi":"10.1257/jel.20180892","locator":"§8 Concluding Remarks, pp.49–51 of June19,2018 revision","evidence_summary":"The conclusion explicitly identifies weaker common-prior assumptions, sequential participation, outside resale, and limited adoption of long-term contracts as research questions."},"formalization_note":"This is an empirical explanation of a stated adoption puzzle. The contract menu and estimands are proposed, and beliefs alone are not assumed to explain all adoption."},"statusQualification":"Published question or research agenda verified at the cited locator. The mathematical specification is adapted for this catalogue and includes additional assumptions; source publication does not establish first-ever priority or certify this exact model remains unsolved. This is an empirical explanation of a stated adoption puzzle. The contract menu and estimands are proposed, and beliefs alone are not assumed to explain all adoption.","statusReviewedAt":"2026-10-08"}
% FIRST_POSTED: 2026-10-08
% ENTRY_KIND: statement_only
% !TeX program = lualatex
\documentclass[11pt]{article}
\usepackage{fontspec}
\usepackage[a4paper,margin=25mm]{geometry}
\usepackage{amsmath,amssymb,mathtools}
\usepackage{xurl}
\usepackage[hidelinks]{hyperref}
\newcommand{\sref}[2]{\hyperlink{source:#1:#2}{[S#2]}}
\setlength{\emergencystretch}{3em}
\begin{document}
% problemId:30007575
\section*{Long-term contracts versus spot markets}\label{30007575}
\subsection{Definitions and mathematical statement}
Consider repeated buyer–seller relationships over \(T\) periods, with observed demand states \(s_t\), private valuations, and measured enforcement and commitment costs. Randomly offer otherwise comparable groups either a prespecified long-term contract menu \(Z=1\) or a sequence of spot-market opportunities \(Z=0\). Let \(W_i(Z)\) be relationship \(i\)'s discounted realized total surplus net of implementation costs, including the buyer's value minus resource cost, and let \(D_i(Z)\) indicate adoption of the offered long-term contract. Elicit buyer and seller forecasts before assignment and define belief divergence \(H_i\) by a stated distance between their predictive distributions. Principal estimands are
\[
 \Delta(h,c)=E[W_i(1)-W_i(0)\mid H_i=h,C_i=c],
 \qquad E[D_i(1)\mid H_i=h,C_i=c],
\]
where \(C_i\) summarizes prespecified enforcement conditions. Identify contract-offer effects using randomization, distinguish adoption effects under explicit instrumental-variable assumptions, and estimate whether belief disagreement explains limited adoption despite potential surplus gains. Use relationship-level assignment, report noncompliance and exit, and include common uncertainty and measurement error in forecasts. Transporting results across platforms requires overlapping contract and state support. The survey states the puzzle of scarce long-term arrangements in dynamic markets; the quantitative estimands and experimental protocol above are an adaptation, rather than a universal theorem about superiority of long-term contracting.

\par\textbf{Formulation and scope.} This is an empirical explanation of a stated adoption puzzle. The contract menu and estimands are proposed, and beliefs alone are not assumed to explain all adoption.

\par\textbf{Open-question provenance.} An explicit question or research agenda was verified at the source locator. This does not by itself certify that every later version remains unresolved \sref{30007575}{1}.

\par\textbf{Historical priority.} This is the earliest source in which the relevant agenda was verified during this audit; absolute priority has not been established.

\subsection{Short English statement}
Identify when offered long-term contracts improve net relationship surplus relative to spot trade and whether measured belief disagreement helps explain limited adoption.

\subsection{Sources}
\begin{itemize}\raggedright
\item[\textnormal{[S1]}]\hypertarget{source:30007575:1}{} Dirk Bergemann, Juuso Välimäki. 2018. Dynamic Mechanism Design: An Introduction. §8 Concluding Remarks, pp.49–51 of June19,2018 revision.
\url{https://cowles.yale.edu/sites/default/files/2022-09/d2102-r.pdf}.
DOI: 10.1257/jel.20180892.
{\small\textit{Earliest verified question/agenda reference; historical priority qualified below. The conclusion explicitly identifies weaker common-prior assumptions, sequential participation, outside resale, and limited adoption of long-term contracts as research questions.}}
\end{itemize}
\end{document}
